% Pollis New Methods: k-nearest neighbours (propensity score matching) — draft 1 (14 October 2026)
% Compile with: latexmk -pdf nm_matching.tex && latexmk -c
% Shared look for the Pollis New Methods documents (same fonts and colours as the brochure and the workflow handouts).
\documentclass[10pt,a4paper]{article}

\usepackage[a4paper,left=18mm,right=18mm,top=15mm,bottom=17mm,footskip=8mm]{geometry}
\usepackage[T1]{fontenc}
\usepackage[default,semibold]{sourcesanspro}
\usepackage{textcomp}
\usepackage[italic]{mathastext}
\usepackage{amsmath}
\usepackage{microtype}
\usepackage[table]{xcolor}
\usepackage{tikz}
\usepackage[most]{tcolorbox}
\usepackage{fontawesome5}
\usepackage{tabularx,booktabs}
\usepackage{enumitem}
\usepackage{fancyhdr}
\usepackage[hidelinks]{hyperref}

% ---------------------------------------------------------------- colours (as in the brochure)
\definecolor{navy}{HTML}{1C3A5E}
\definecolor{pblue}{HTML}{3B78B5}
\definecolor{amber}{HTML}{F2B33D}
\definecolor{deepamber}{HTML}{D9902A}
\definecolor{pale}{HTML}{EEF3F9}
\definecolor{ink}{HTML}{2A2F36}
\definecolor{logonavy}{HTML}{173D63}
\definecolor{logodark}{HTML}{C98A1E}
\definecolor{logoamber}{HTML}{F2B446}
\definecolor{logocream}{HTML}{FEEFD2}
\definecolor{soft}{HTML}{5E6773}
\definecolor{agree}{HTML}{2F7D4F}
\definecolor{differ}{HTML}{B23B3B}
\color{ink}

\setlength{\parindent}{0pt}
\setlength{\parskip}{4pt}
\setlist{nosep,leftmargin=*,itemsep=2pt}
\renewcommand{\arraystretch}{1.15}

\pagestyle{fancy}
\fancyhf{}
\renewcommand{\headrulewidth}{0pt}
\fancyfoot[L]{\footnotesize\color{soft}\textbf{Pollis}\enspace New Methods\enspace·\enspace\nmtheme}
\fancyfoot[R]{\footnotesize\color{soft}\thepage}
\newcommand{\nmtheme}{}
\newcommand{\nmdate}{11 October 2026}

% The Pollis logo, redrawn so it prints sharply at any size
\newcommand{\pollislogo}[1]{%
  \begin{tikzpicture}[x=#1/512,y=-#1/512,baseline=(current bounding box.center)]
    \fill[logonavy] (0,0) rectangle (512,512);
    \fill[logodark] (36,422) -- (104,256) -- (172,422) -- cycle;
    \fill[logoamber] (104,422) -- (203,180) -- (302,422) -- cycle;
    \fill[logocream] (203,422) -- (339,89) -- (476,422) -- cycle;
  \end{tikzpicture}}

% Top of the first page: logo, name; then the kicker, title and a one-line promise
\newcommand{\nmhead}[3]{%
  \renewcommand{\nmtheme}{#1}%
  \thispagestyle{fancy}%
  \begin{minipage}[c]{13mm}\pollislogo{11mm}\end{minipage}%
  \begin{minipage}[c]{\dimexpr\linewidth-13mm\relax}
    {\color{navy}\Large\bfseries Pollis}\enspace{\color{soft}New Methods}\hfill
    {\color{deepamber}\small\bfseries\MakeUppercase{#1}}
  \end{minipage}\par\vspace{7mm}
  {\color{navy}\fontsize{22}{26}\selectfont\bfseries #2}\par\vspace{4pt}
  {\color{soft}\large #3}\par\vspace{5mm}}

\newcommand{\nmsection}[1]{\par\vspace{5pt}{\color{navy}\large\bfseries #1}\par\vspace{1pt}}
\newcommand{\nmsub}[1]{\par\vspace{2pt}{\color{deepamber}\bfseries #1}\par}
\newcommand{\gt}{\,{\color{deepamber}\bfseries ›}\,}
\newcommand{\menupath}[1]{{\bfseries\color{navy}#1}}
\newcommand{\fn}[1]{\texttt{\small #1}}
% Citations: \nmcite{key}{Breiman (2001)} links to the entry \nmref{key}{...} in the References
\newcommand{\nmcite}[2]{\hyperlink{ref:#1}{\color{pblue}#2}}
\newcommand{\nmref}[2]{\raisebox{1.1em}[0pt][0pt]{\hypertarget{ref:#1}{}}#2\par}
\newcommand{\nmdoi}[1]{\href{https://doi.org/#1}{\color{pblue}doi:#1}}
% Where a notebook is in Pollis: its title in Browse Methods, and its file
\newcommand{\browse}[2]{\menupath{Compose \gt New Methods \gt Browse Methods}: \emph{#1} (\fn{#2.ipynb})}

% Key facts at the top: source, missing in Julia, notebook, result
\newcommand{\factsheet}[4]{%
  \begin{tcolorbox}[enhanced, colback=pale, colframe=pale, boxrule=0pt, sharp corners,
    left=7pt, right=7pt, top=5pt, bottom=5pt, fontupper=\small]
  \begin{tabularx}{\linewidth}{@{}>{\leavevmode\bfseries\color{navy}}p{27mm}>{\leavevmode}X@{}}
  Tested against & #1\\
  Missing in Julia & #2\\
  Notebook & #3\\
  Result & #4
  \end{tabularx}
  \end{tcolorbox}}

% Boxes
\newtcolorbox{pbox}[2][]{enhanced, colback=pale, colframe=amber, boxrule=0pt,
  leftrule=3pt, sharp corners, left=7pt, right=7pt, top=6pt, bottom=6pt,
  fontupper=\small\raggedright, before upper={{\color{navy}\normalsize\bfseries #2}\par\smallskip}, #1}
\newtcolorbox{nbox}[2][]{enhanced, colback=navy, colframe=navy, boxrule=0pt,
  sharp corners, left=8pt, right=8pt, top=7pt, bottom=7pt, coltext=white,
  fontupper=\small\raggedright, before upper={{\color{amber}\normalsize\bfseries #2}\par\smallskip}, #1}

% Value tables: published beside Julia
\newenvironment{valuetable}[1]{%
  \par\small\tabularx{\linewidth}{#1}\toprule}{\bottomrule\endtabularx\par}
\newcommand{\ok}{{\color{agree}\faCheck}}
\newcommand{\no}{{\color{differ}\faTimes}}

% End of the document: credit and licence
\newcommand{\nmcredit}{%
  \par\vfill
  \begin{minipage}[b]{0.78\linewidth}\footnotesize\color{soft}
  Methods written and tested with Claude Code (Anthropic) in Pollis, from the published papers and their data.\par
  \textcopyright\ 2026 Antonio Saragga Seabra\enspace·\enspace Licensed under \href{https://creativecommons.org/licenses/by-sa/4.0/}{CC BY-SA 4.0}\enspace·\enspace Draft, \nmdate
  \end{minipage}\hfill
  \begin{minipage}[b]{0.18\linewidth}\raggedleft\pollislogo{9mm}\par\vspace{2pt}{\footnotesize contact@pollis.org}\end{minipage}}


\begin{document}

\nmhead{K-Nearest Neighbours}
  {Propensity score matching: Dehejia and Wahba, and its replications}
  {Can a nonexperimental study find the effect a randomised experiment found? Dehejia and Wahba (1999), Smith and Todd (2005) and Dehejia (2005), reproduced.}

\factsheet
  {\nmcite{dw1999}{Dehejia and Wahba (1999)}, \emph{Journal of the American Statistical Association}, Tables 1 and 3;
   \nmcite{st2005}{Smith and Todd (2005)}, \emph{Journal of Econometrics}, Tables 3, 5 and 7; and
   \nmcite{d2005}{Dehejia (2005)}, its reply, Tables 2--4. The data are LaLonde's NSW men and the PSID and CPS comparison
   groups from Dehejia's data page (read into memory and not saved). The matches were run again in R with
   \fn{MatchIt} and \fn{Matching}, and the balance checked with \fn{cobalt}.}
  {the logit propensity score, common support by the treated range or by Smith and Todd's density trimming,
   nearest-neighbour matching with replacement, local linear matching, the effect on the treated and the balance of the
   covariates. Julia has no counterpart of MatchIt, Matching or cobalt.}
  {\browse{Propensity score matching: Dehejia and Wahba, and its replications}{PropensityMatching}}
  {123 of 145 printed values to every digit, the other 22 in the last. Smith and Todd's PSID logit to every digit.}

\nmsection{The problem}
\nmcite{lalonde1986}{LaLonde (1986)} took the National Supported Work experiment, whose randomised controls give the
effect of the training on the earnings of the men treated, about \$1,794, and asked whether nonexperimental methods could
find it from survey data instead: the PSID and CPS samples in place of the controls. They could not; the estimates ranged
widely and depended on the specification. \nmcite{dw1999}{Dehejia and Wahba (1999)} matched each treated man to the
comparison men with the nearest estimated probability of treatment, the propensity score, and came close to the
experiment. \nmcite{st2005}{Smith and Todd (2005)} replied that this works only in the subsample Dehejia and Wahba chose,
and only for some specifications; \nmcite{d2005}{Dehejia (2005)} replied that the specification must be chosen for each
sample by the balance of the covariates. The three papers together are a test of matching, and of how much an estimate
depends on choices a paper does not print.

\nmsection{The method}
The propensity score $p(x) = P(d = 1 \mid x)$ is fitted by a logit (\fn{logit}), with earnings in thousands of dollars.
Comparison units whose score lies outside the range of the treated are discarded (\fn{commonsupport}), or, as Smith and
Todd do, the treated whose score density among the comparison units is below the 2 per cent point are dropped
(\fn{trimmedsupport}, a quartic kernel with Silverman's bandwidth). Each treated unit is matched, with replacement, to the
$k$ comparison units with the nearest score (\fn{nnmatch}), on $p$ or on the log-odds $\ln(p/(1-p))$ (\fn{logodds}),
and the effect on the treated is the mean difference in 1978 earnings (\fn{att}). When several comparison units are
exactly as near, \fn{nnmatch} takes the first in the data; \fn{tierange} gives the smallest and largest effect over every
way of breaking those ties. \fn{locallinear} is Smith and Todd's local linear matching (their eq.\,11), \fn{matchweights}
turns matches into weights and \fn{balance} gives the standardised differences of the covariates, as \fn{cobalt} does.

\nmsection{What the tests show}
\begin{valuetable}{@{}lXl@{}}
\textbf{Test} & \textbf{What is compared} & \textbf{Agreement}\\\midrule
DW Table 1 & means of 8 covariates, 10 samples: 77 values & 62 to every digit, 15 in the last\\
DW Table 3 & col.\,(1), the unadjusted differences: 7 values & 6 to every digit, 1 in the last\\
ST Table 3 & DW 2002 logits, CPS: 15 coefficients, 15 SEs & 29 to every digit, 1 in the last\\
ST Table 3 & DW 2002 logit, PSID: 16 coefficients & every one\\
ST Tables 5, 7 & bias of the mean difference and regression: 6 values & 5 to every digit, 1 in the last\\
ST Table 5 & CPS local linear matching, bandwidths 1 and 4 & both in the last digit\\
Dehejia Tables 3, 4 & nearest-neighbour effects: 7 values & 5 to every digit, 2 in the last\\
\end{valuetable}

The last-digit differences are mostly truncation: Dehejia and Wahba print 14016.8 as 14016 and 33.2252 as 33.22. The
others, side by side (Julia against the paper):
\begin{itemize}
  \item Smith and Todd's CPS intercept: $-39.832651$ vs $-39.8326$.
  \item Smith and Todd's CPS local linear bias: $-87.0$ vs $-88$ (bandwidth 1) and $-67.94$ vs $-67$ (bandwidth 4).
  \item Dehejia's ST specifications, nearest neighbour: 1,282.49 vs 1,283 (CPS) and 1,990.47 vs 1,991 (PSID).
\end{itemize}

\nmsub{Against R}
With Julia's scores passed in, every match \fn{MatchIt} makes is one of the exactly nearest comparison units, for all 185
or 297 treated in all 8 cases, and its effects equal Julia's where the tie choices coincide: 1,869.07, 1,990.47 and
862.72. \fn{Matching} averages over near-ties within an internal tolerance that \fn{distance.tolerance} does not switch
off, so it gives 1,829.82 against 1,869.07 for PSID/DW; on CPS/Lalonde it gives 1,858.87, Julia's average over ties.
\fn{cobalt}'s \fn{bal.tab} with Julia's weights gives the same 16 standardised differences to 4 decimals (age $-1.2627$
before matching and $0.0431$ after, \fn{re75} $-5.4458$ and $-0.1997$).

\nmsection{Three things the papers do not say}
\nmsub{Ties decide the CPS estimate}
In Dehejia's CPS/DW specification 22 treated men have two or more comparison men at exactly the same score, in sets of up
to 9. Taking the first in the data gives 1,762; over every way of breaking the ties the effect runs from 1,055 to 1,964,
and their average is 1,622. The printed 1,589 is inside that range, but no rule a reader could follow from the paper gives
it, so it is not among the tests. Smith and Todd's CPS single-neighbour bias, 407, is likewise inside a tie range of 41 to
828 (Julia's first-of-ties gives 535).

\nmsub{A printed mean difference is the 1975 gap}
Smith and Todd's Table 7 gives the PSID mean difference as $-17{,}796$ (846). That is the gap in 1975 earnings, the value
of their own Table 8; the gap in 1978 earnings, the outcome, is $-16{,}999$, as in their Table 5.

\nmsub{A corrected mean}
Dehejia and Wahba's Table 1 gives 0.18 for the share of Hispanics in PSID-3; the data give 0.117. The other 76 means
agree, so the value is left untested; it looks like a correction to the data after the table was printed.

\nmsection{Still open}
\begin{itemize}
  \item \textbf{Dehejia and Wahba's own matches.} Their Table 3, columns (2)--(8), does not reproduce: PSID-1 gives 560
    against 1,691. Smith and Todd (footnote 40) note that the sample on which the score is fitted moves these estimates by
    \$200--300; the Dehejia--Wahba specifications are tested through Smith and Todd's logits and Dehejia (2005) instead.
  \item \textbf{Smith and Todd's other estimators.} Their PSID single-neighbour bias, 361, is what matching on $p$ gives;
    on the log-odds they say they use, it is 319. The ten-neighbour and trimmed estimators come within about \$100 (CPS
    ten neighbours $-76$ vs $-5$, inside a tie range of $-176$ to 14), and the PSID local linear bias does not reproduce
    ($-160$ and 124 vs $-122$ and 143). The tie rule and the trimming bandwidth are not stated.
  \item \textbf{Smith and Todd's PSID standard errors} agree only to about three significant figures (intercept 2.0623 vs
    2.0629), and are not tested.
  \item \textbf{The Lalonde and ST samples.} Dehejia's cross-specifications on LaLonde's sample need 1974 earnings, which
    the public files do not have for it, and Smith and Todd's early random assignment sample cannot be rebuilt from them.
\end{itemize}

\nmsection{Further considerations}
\begin{itemize}
  \item \textbf{Report the tie range.} Where a score takes few values, as here, an estimate can depend on the order of the
    data; \fn{tierange} shows by how much, and is worth reporting beside any nearest-neighbour estimate.
  \item \textbf{Other matching estimators.} Kernel matching, calipers and Abadie and Imbens's bias-corrected standard errors
    are in \fn{Matching}; each would need its own published table before it becomes a test here.
  \item \textbf{Upstream.} Julia has no counterpart of MatchIt or cobalt; the three papers are a ready-made test for one.
\end{itemize}

\nmsection{References}
{\footnotesize
\nmref{dw1999}{Dehejia, R.\,H. and Wahba, S. (1999). Causal effects in nonexperimental studies: reevaluating the evaluation of training programs. \emph{Journal of the American Statistical Association} 94(448), 1053--1062. \nmdoi{10.1080/01621459.1999.10473858}}
\nmref{dw2002}{Dehejia, R.\,H. and Wahba, S. (2002). Propensity score-matching methods for nonexperimental causal studies. \emph{Review of Economics and Statistics} 84(1), 151--161. \nmdoi{10.1162/003465302317331982}}
\nmref{st2005}{Smith, J.\,A. and Todd, P.\,E. (2005). Does matching overcome LaLonde's critique of nonexperimental estimators? \emph{Journal of Econometrics} 125(1--2), 305--353. \nmdoi{10.1016/j.jeconom.2004.04.011}}
\nmref{d2005}{Dehejia, R. (2005). Practical propensity score matching: a reply to Smith and Todd. \emph{Journal of Econometrics} 125(1--2), 355--364.\newline \nmdoi{10.1016/j.jeconom.2004.04.012}}
\nmref{lalonde1986}{LaLonde, R.\,J. (1986). Evaluating the econometric evaluations of training programs with experimental data. \emph{American Economic Review} 76(4), 604--620.}
\nmref{data}{Data: Dehejia's NSW, PSID and CPS files, from Dehejia's data page at the NBER, read into memory and not saved, since they are not offered for redistribution with Pollis.}
}

\nmcredit
\end{document}
